% Shared section: the Fabre (2025) survey experiment
\subsection{A survey experiment on question wording and scale}\label{subsec:fabre}

Gallup and WVS differ in two respects that could explain why they disagree: the question (the Cantril ladder on a 0--10 scale vs.\ life satisfaction on a 1--10 scale) and everything else---sampling frames, interview modes, questionnaire context, translations, and survey years. To separate the two, I embedded a randomized experiment in an original survey \citep{fabre_public_2025}. The survey was conducted online between April 15 and July 3, 2025, by Bilendi (Kantar in Saudi Arabia) on \nFabre respondents in \nFabreCountries high-income countries: the United States (3,000), Japan (2,000), Saudi Arabia (1,000), Germany (1,048), the United Kingdom (826), France (798), Italy (756), Spain (603), Poland (500) and Switzerland (469). (The survey also covered Russia, where the well-being question was not asked.) Samples were stratified with quotas on gender, age, income, education, region and urbanicity (race in the U.S.; citizenship in Saudi Arabia), and results are reweighted to match population frequencies along these dimensions (weights trimmed between 0.25 and 4). Respondents failing an attention test or completing the questionnaire in less than 6 minutes were excluded. The project was approved by the CIRED institutional review board (IRB-CIRED-2025-2) and pre-registered at the Open Science Framework (\href{https://osf.io/7mzn4}{osf.io/7mzn4}).

Toward the end of the questionnaire, after questions on global policies, each respondent was randomly assigned to one of four versions of the life-evaluation question:
\begin{itemize}
    \item \textit{Ladder 0--10} (Gallup's original): the Cantril ladder question quoted in Section~\ref{subsec:gallup}, with steps numbered from 0 (\textit{Worst possible}) to 10 (\textit{Best possible});
    \item \textit{Ladder 1--10}: same wording, with steps numbered from 1 to 10;
    \item \textit{Satisfaction 0--10}: the WVS question quoted in Section~\ref{subsec:wvs}, with answers from 0 (\textit{Completely dissatisfied}) to 10 (\textit{Completely satisfied});
    \item \textit{Satisfaction 1--10} (WVS's original): same wording, with answers from 1 to 10.
\end{itemize}
This $2 \times 2$ design identifies the effect of wording (ladder vs.\ satisfaction) and scale (0--10 vs.\ 1--10), and their interaction, holding the sample, mode, context and time constant. Within each country, the four branches are balanced (each receives between 21\% and 30\% of respondents).
